
\documentclass{report}

\input{preamble}
\input{macros}
\input{letterfonts}

\title{\Huge{Fisica}\\Sistemi Informativi Aziendali}
\author{\huge{Caberlotto Giovanni}}
\date{}

\begin{document}

\maketitle
\newpage% or \cleardoublepage
% \pdfbookmark[<level>]{<title>}{<dest>}
\pdfbookmark[section]{\contentsname}{toc}
\tableofcontents
\pagebreak

\chapter{Sistemi di misure grandezze fondamentali e grandezze scalari e vettoriali} % (fold)
\label{chap:Sistemi di misure e grandezze fondamentali e grandezze scalari e vettoriali}

\section{Sistema metrico decimale} % (fold)
\label{sec:Sistema metrico decimale}

% section Sistema metrico decimale (end)

\section{Sistema Internazionale} % (fold)
\label{sec:Sistema Internazionale}
Il Sistema Internazionale è un trattato in cui, tra le tante cose, sono state stabilite 7 grandezze fisiche dette grandezze fondamentali con annesse unità di misura, in funzione delle quali vengono definite tutte le altre grandezze, dette per l'appunto grandezze derivate.

\subsection{Grandezze fondamentali del Sistema Internazionale} % (fold)
\label{sub:Grandezze fondamentali del Sistema Internazionale}

Come già anticipato il S.I. prevede l'utilizzo di 7 unità di misura per 7 grandezze fisiche fondamentali del Sistema Internazionale. Per ogni unità di misura è stato inoltre stabilito il simbolo da utilizzare universalmente, le regole sull'utilizzo dei vari simboli e sulla composizione delle unità di misura derivate dalle fondamentali.
\\
\begin{table}[h]
  \caption{Grandezze Fondamentali}\label{tab:}
  \begin{center}
    \begin{tabular}[c]{|l|l|l|}
      \hline
      Grandezze Fondamentali & Unità di misura & Simbolo \\
      \hline
      Lunghezza & metro & $m$ \\
      \hline
      Tempo & secondo & $s$ \\
      \hline 
      Massa & kilogrammo & $kg$ \\
      \hline
      Intensità di corrente & ampere & $A$ \\
      \hline 
      Temperatura & kelvin & $K$ \\
      \hline 
      Quantità di materia & mole & $mol$ \\
      \hline
      Intensità luminosa & candela & $cd$ \\
      \hline
    \end{tabular}
  \end{center}
\end{table}

% subsection Grandezze fondamentali del Sistema Internazionale (end)

\subsection{Grandezze derivate del Sistema Internazionale} % (fold)
\label{sub:Grandezze derivate del Sistema Internazionale}
Le grandezze derivate del Sistema Internazionale con le rispettive unità di misura si ottengono moltiplicando o dividendo tra loro due o più grandezze fondamentali. Ad esempio:
\begin{itemize}
  \item{La velocità è una grandezza derivata definita come il rapporto tra lunghezza e tempo, e si misura in $\frac{m}{s}$;}
  \item{La forza peso è anch'essa una grandezza derivata del S.I. la cui unità di misura è data dal $kg\frac{m}{2s}$ e viene indicata brevemente con $N$ (newton).}
\end{itemize}

% subsection Grandezze derivate del Sistema Internazionale (end)

\subsection{Prefissi per le unità del Sistema Internazionale} % (fold)
\label{sub:Prefissi per le unità del Sistema Internazionale}

Quando si ha a che fare con misure di grandezze molto grandi o molte piccole, per comodità, possiamo usare i seguenti prefissi che rappresentano un fattore di conversione dato da una certa potenza di 10.
\\
\begin{table}[h]
  \caption{Prefissi}\label{tab:}
  \begin{center}
    \begin{tabular}[c]{|l|l|l|}
      \hline
      Prefisso & Fattore di conversione & Simbolo \\
      \hline
      yotta   & $10^{24}$  & $Y$ \\
      \hline
      zeta    & $10^{21}$  & $Z$ \\
      \hline
      exa     & $10^{18}$  & $E$ \\
      \hline 
      peta    & $10^{15}$  & $P$ \\
      \hline 
      tera    & $10^{12}$  & $T$ \\
      \hline 
      giga    & $10^9$     & $G$ \\
      \hline 
      mega    & $10^6$     & $M$ \\
      \hline 
      kilo    & $10^3$     & $k$ \\
      \hline 
      etto    & $10^2$     & $h$ \\
      \hline 
      deca    & $10^1$     & $da$ \\
      \hline
      deci    & $10^{-1}$  & $d$ \\
      \hline
      centi   & $10^{-2}$  & $c$ \\
      \hline 
      milli   & $10^{-3}$  & $m$ \\
      \hline 
      micro   & $10^{-6}$  & $\mu$ \\ 
      \hline
      nano    & $10^{-9}$  & $n$ \\
      \hline
      pico    & $10^{-12}$ & $p$ \\
      \hline
      femto   & $10^{-15}$ & $f$ \\
      \hline
      atto    & $10^{-18}$ & $a$ \\
      \hline
      zepto   & $10^{-21}$ & $z$ \\
      \hline
      yocto   & $10^{-24}$ & $y$ \\
      \hline
    \end{tabular}
  \end{center}
\end{table}
\\
Associare un prefisso ad un'unità di misura significa quindi moltiplicarla per il relativo fattore, in modo da ottenere il corrispondente multiplo o sottomultiplo dell'unità di misura. Ad esempio:
\\
\\
$7.315$ nanosecondi $= 7.315 * 10^{-9}$
% subsection Prefissi per le unità del Sistema Internazionale (end)

% section Sistema Internazionale (end)

\section{Analisi dimensionale} % (fold)
\label{sec:Analisi dimensionale}

% section Analisi dimensionale (end)

% chapter Sistemi di misure e grandezze fondamentali e grandezze scalari e vettoriali (end)

\chapter{Cinematica} % (fold)
\label{chap:Cinematica}

% chapter Cinematica (end)

\end{document}
